\documentclass[11pt,a4paper]{article}
\usepackage[T1]{fontenc}
\usepackage[utf8]{inputenc}
\usepackage{lmodern}
\usepackage{microtype}
\usepackage[a4paper,margin=0pt]{geometry}
\usepackage[table]{xcolor}
\usepackage{graphicx}
\usepackage{tabularx}
\usepackage{ragged2e}
\usepackage{hyperref}

\definecolor{Ink}{HTML}{182433}
\definecolor{Slate}{HTML}{536273}
\definecolor{Muted}{HTML}{788696}
\definecolor{Blue}{HTML}{2A7AE2}
\definecolor{Purple}{HTML}{7C3AED}
\definecolor{PaleBlue}{HTML}{F2F7FD}
\definecolor{PaleGray}{HTML}{F6F8FA}
\definecolor{Rule}{HTML}{DDE4EA}

\hypersetup{
  pdftitle={AI and Mathematical Proof: From a Sparse 2025 Baseline to the 2026 Take-off},
  pdfauthor={SuperLattice},
  pdfsubject={Automated theorem proving and AI-assisted mathematical discovery}
}
\setlength{\parindent}{0pt}
\setlength{\parskip}{0pt}
\setlength{\fboxsep}{0pt}
\renewcommand{\familydefault}{\sfdefault}
\newcolumntype{L}[1]{>{\RaggedRight\arraybackslash}p{#1}}

\newcommand{\metric}[2]{%
  \textbf{\color{Ink}#1} & \color{Slate}#2\\[1.6mm]
}

\begin{document}
\begin{titlepage}
\thispagestyle{empty}

% Restrained masthead: one baseline, one rule, one type family.
\noindent\colorbox{Ink}{\parbox[c][24mm][c]{\paperwidth}{%
  \hspace*{19mm}%
  {\fontsize{17}{20}\selectfont\bfseries\color{white}SuperLattice}%
  \hfill
  {\fontsize{9.5}{12}\selectfont\color{white!72}RESEARCH REPORT}%
  \hspace*{19mm}}}

\vspace{18mm}
\begin{minipage}{\paperwidth}
\hspace*{22mm}\begin{minipage}{167mm}
  {\fontsize{9.5}{12}\selectfont\bfseries\color{Blue}\MakeUppercase{Research project}\hspace{2mm}%
   \color{Muted}\textnormal{Evidence through 21 July 2026}}\par
  \vspace{6mm}

  {\RaggedRight\fontsize{31}{34}\selectfont\bfseries\color{Ink}%
   AI and mathematical proof\par progress}\par
  \vspace{5mm}

  {\RaggedRight\fontsize{17}{21}\selectfont\color{Slate}%
   From a sparse 2025 baseline to the 2026 take-off}\par
  \vspace{5mm}

  {\RaggedRight\fontsize{11}{16}\selectfont\color{Muted}%
   An exploratory record of audited AI contributions to automated and semi-automated mathematical proofs, constructions, bounds and counterexamples.}\par

  \vspace{10mm}
  {\color{Rule}\rule{\linewidth}{0.6pt}}\par
  \vspace{7mm}

  \begin{minipage}[t]{103mm}
    {\fontsize{9}{12.5}\selectfont
    \begin{tabularx}{\linewidth}{@{}L{31mm}X@{}}
      \metric{Evidence window}{1 January 2025 to 21 July 2026.}
      \metric{Research ledger}{Twenty entries first reported in 2025 plus 44 entries first reported in 2026.}
      \metric{Counting rule}{A month receives an entry only for a new public mathematical conclusion, best-known bound, counterexample, or bounded construction portfolio with a documented AI contribution and a checkable source.}
      \metric{Zero-month rule}{Zero means that this audit found no public item meeting the threshold. It does not prove that no private, unreported, or differently classified AI-assisted work occurred.}
    \end{tabularx}}
  \end{minipage}\hfill
  \begin{minipage}[t]{52mm}
    \vspace{1mm}
    \colorbox{PaleBlue}{\parbox{45mm}{%
      \vspace{4mm}\hspace{3mm}%
      \begin{minipage}{37mm}
      {\fontsize{8.6}{12}\selectfont\color{Ink}
      \textbf{AI analysis disclaimer}\par\vspace{2mm}
      \color{Slate}This is an AI-assisted analysis, not an expert-reviewed census. The evidence collection was last run through \textbf{21 July 2026} and contains \textbf{64 audited entries}. The supplied source bundle does not record which model produced the baseline collection. The dashboard and hybrid-scoring rerun were prepared with OpenAI Codex, based on GPT-5. Treat classifications as reproducible proxies and check substantive claims against the cited works.}
      \end{minipage}\hspace{3mm}\vspace{4mm}}}
  \end{minipage}

  \vspace{8mm}
  \colorbox{PaleGray}{\parbox{\dimexpr\linewidth-8mm\relax}{%
    \vspace{4mm}\hspace{4mm}\begin{minipage}{\dimexpr\linewidth-16mm\relax}
    {\fontsize{9.5}{13.5}\selectfont\color{Slate}%
    \textbf{\color{Ink}Central finding.} January--March 2025 form a flat audited baseline. Isolated discoveries begin in April and May, June--July are flat again, and a broader August--December ramp precedes the much larger portfolio-production surge of spring 2026.}
    \end{minipage}\hspace{4mm}\vspace{4mm}}}

  \vfill
  \vspace{8mm}
  {\color{Rule}\rule{\linewidth}{0.6pt}}\par\vspace{3mm}
  {\fontsize{8.5}{11}\selectfont\color{Muted}64 audited entries\hfill January 2025--21 July 2026}
\end{minipage}
\end{minipage}

\end{titlepage}
\end{document}
